\documentclass[11pt,a4paper]{article}

% ==============================================================================
% PACKAGES AND CONFIGURATION
% ==============================================================================
\usepackage[utf8]{inputenc}
\usepackage[margin=1in]{geometry}
\usepackage{amsmath,amssymb,amsfonts,amsthm}
\usepackage{graphicx}
\usepackage{booktabs}
\usepackage{multirow}
\usepackage{array}
\usepackage{siunitx}
\usepackage{url}
\usepackage{hyperref}
\usepackage{float}
\usepackage{caption}
\usepackage{subcaption}
\usepackage{algorithm}
\usepackage{algorithmic}
\usepackage{microtype}
\usepackage{setspace}
\usepackage{authblk}

\hypersetup{
    colorlinks=true,
    linkcolor=blue!70!black,
    citecolor=blue!70!black,
    urlcolor=blue!70!black
}

\onehalfspacing

% ==============================================================================
% TITLE AND AUTHOR METADATA
% ==============================================================================
\title{\textbf{Bridging Multi-Horizon Retail Demand Forecasting and Inventory Decisions via Rolling Forecast-Conditional Conformal Uncertainty Calibration}}

\author[1,*]{R. Mangaiyarkarasi}
\author[1]{Research Collaborators}
\affil[1]{Department of Computer Science and Engineering}
\affil[*]{Corresponding author: \texttt{author@institution.edu}}

\date{September 2026}

\begin{document}

\maketitle

% ==============================================================================
% ABSTRACT
% ==============================================================================
\begin{abstract}
\noindent Modern retail supply chain planning relies on demand forecasting models to govern inventory replenishment schedules. However, enterprise forecasting pipelines predominantly optimize symmetric point-prediction errors, failing to communicate predictive uncertainty or mitigate asymmetric stockout penalties. Furthermore, temporal nonstationarity and scale shifts routinely violate the exchangeability assumptions underlying probabilistic calibration. This study presents an end-to-end empirical framework bridging leakage-free multi-horizon forecasting, rolling forecast-conditional conformal prediction, and continuous-review inventory simulation across 500 store-item time series. Extending our preliminary conference baseline (\emph{E-IDSS}), we implement seven decoupled direct Extreme Gradient Boosting (XGBoost) regressors enforcing a decision-time minimum lag constraint ($\tau \ge 7$ days), achieving stable validation accuracy ($\text{MAE} \approx 6.04$) while eliminating recursive error compounding and passing 10/10 automated leakage tests. On an untouched 365-day test holdout, rolling conformal prediction experiences empirical under-coverage ($17.90\%-31.98\%$ at nominal $90\%$ confidence) driven by a $3.55\times$ to $4.30\times$ expansion in residual standard deviation and seasonal drift ($p < 10^{-42}$). Relative error scaling significantly improves interval efficiency, lowering mean Winkler loss to $409.66$ ($p = 1.44 \times 10^{-28}$) and increasing high-demand coverage by $+8.80$ percentage points. When translated into dynamic replenishment buffers ($ROP = \hat{D} + SS$), uncertainty-aware replenishment increases the portfolio demand-weighted fill rate by $+10.64$ percentage points ($46.87\% \to 57.51\%$) at lead time $L=1$ ($p < 10^{-75}$) and cuts lost sales volume by $20.0\%-29.6\%$ ($1.14$ million units saved). Economic analysis reveals a $-10.83\%$ total operating cost reduction at $L=1$ ($-\$790,834$, $p = 6.65 \times 10^{-55}$) and maps the multi-day trade-off frontier where conformal replenishment is cost-preferred whenever shortage penalties exceed $\$2.00/\text{unit}$.
\end{abstract}

\textbf{Keywords:} Retail Demand Forecasting, Multi-Horizon Forecasting, Conformal Prediction, Uncertainty Quantification, Distribution Shift, Inventory Decision Support, XGBoost, Supply Chain Analytics.

\vspace{0.5cm}
\hrule
\vspace{0.5cm}

% ==============================================================================
% SECTION 1: INTRODUCTION
% ==============================================================================
\section{Introduction}
\label{sec:introduction}

\subsection{Background and Motivation}
\label{sec:intro_bg}
Retail supply chain management operates in an increasingly volatile, data-intensive environment where accurate demand estimation is fundamental to operational viability \cite{tiwari2025agentic}. Modern multi-channel retail enterprises must coordinate procurement, warehousing, and inventory allocation across hundreds of stores and thousands of stock-keeping units (SKUs). These retail portfolios exhibit substantial demand heterogeneity, characterized by pronounced day-of-week seasonality, periodic promotional surges, holiday spikes, and intermittent low-velocity purchasing patterns \cite{mitra2022comparative, venkiteela2025machine}. Within this operational setting, demand forecasting serves as the cornerstone of Sales and Operations Planning (S\&OP), directly governing procurement schedules, labor allocation, and working capital commitments \cite{mustapha2025forecasting}. Inaccurate demand projections propagate downstream into operational failures: under-forecasting leads to inventory stockouts, immediate lost sales, and customer defection, whereas over-forecasting results in excessive inventory holding costs, working capital lockup, and product obsolescence \cite{spanos2022application, malik2024sales}. While predictive accuracy is a necessary prerequisite for effective planning, operational decisions depend not only on expected demand quantities but also fundamentally on the uncertainty surrounding those predictions.

\subsection{The Problem: Limitations of Point Forecasts}
\label{sec:intro_problem}
In conventional enterprise forecasting workflows, predictive models output a single deterministic point forecast, denoted as $\hat{y}_{t+h}$ for a future horizon $h$ from forecast origin $t$. While point forecasts estimate conditional expected demand, they do not communicate the dispersion, skewness, or tail risk inherent in the forecast error distribution \cite{gneiting2007strictly}. Standard machine learning loss functions, such as Mean Absolute Error (MAE) and Root Mean Squared Error (RMSE), train models under symmetric error penalization, implicitly treating positive and negative forecast errors as equally costly. Yet, in retail replenishment, the financial and operational consequences of estimation errors are fundamentally asymmetric \cite{ban2019big, bertsimas2020predictive}. An under-forecast of ten units during a peak sales window directly triggers stockouts and permanently lost revenue under lost-sales dynamics, whereas an over-forecast of ten units merely incurs a marginal daily holding cost. Consequently, relying exclusively on deterministic point forecasts provides limited information for inventory replenishment decisions that must actively balance asymmetric operational penalties.

\subsection{Limitations in Existing Work}
\label{sec:intro_limitations}
To capture the complex, non-linear dependencies governing retail sales, the forecasting literature has progressed from univariate linear statistical models, such as Autoregressive Integrated Moving Average (ARIMA) and Exponential Smoothing (ETS), to global machine learning algorithms \cite{mustapha2025forecasting, yeesin2024application}. In extensive empirical benchmarks, including the large-scale M5 competition, Gradient Boosted Decision Trees (GBDTs), most notably Extreme Gradient Boosting (XGBoost) and LightGBM, have demonstrated superior predictive accuracy and computational efficiency compared to classical statistical benchmarks and recurrent neural architectures on structured tabular demand series \cite{makridakis2022m5, chen2016xgboost, hobor2026comparative, alghurabli2025ai}. However, in multi-step forecasting settings, standard recursive forecasting strategies, where model predictions are iteratively fed back as lagged predictors, frequently suffer from compounding error propagation and unintended look-ahead feature leakage if feature pipelines access observations within the forecast lead-time window \cite{tiwari2026explainable}. Ensuring that multi-horizon feature pipelines remain strictly decision-time-safe is therefore critical for real-world operational integrity.

To quantify forecast dispersion without imposing restrictive Gaussian assumptions \cite{silver1998inventory}, research in distribution-free uncertainty quantification has increasingly turned to Conformal Prediction (CP) \cite{vovk2005algorithmic, angelopoulos2021gentle}. Conformal prediction provides a mathematically rigorous, model-agnostic framework for constructing valid prediction intervals using empirical nonconformity scores derived from historical calibration residuals \cite{stankeviciute2021conformal, huber2019integrating}. However, standard conformal prediction theory fundamentally relies on data exchangeability, an assumption that is routinely violated in nonstationary retail environments exhibiting macroeconomic growth and variance expansion \cite{barber2023conformal, tibshirani2019conformal}. Furthermore, the operational bridge connecting non-parametric conformal uncertainty margins directly to dynamic safety stock buffers across multi-day lead times ($L \in \{1, 3, 5, 7\}$) in lost-sales inventory simulation remains underexplored \cite{syntetos2016supply, axsaeter2015inventory, snyder2019fundamentals}.

\subsection{Previous Conference Study}
\label{sec:intro_conference}
In our prior baseline study \cite{e_idss_ieee}, we introduced the Explainable Intelligent Decision Support System (E-IDSS), which deployed an optimized single-step XGBoost forecasting model integrated with TreeSHAP (Tree-based Shapley Additive exPlanations) \cite{lundberg2017unified} to provide post-hoc local and global feature attributions for retail sales predictions. The E-IDSS framework addressed the enterprise ``explanation gap,'' demonstrating how transparent feature importances can demystify black-box machine learning models for store managers and reduce subjective forecast override rates \cite{lin2026critical, chaganti2026, narlawar2025explainable, chang2025explainable, madabhushini2025}. Nevertheless, a critical methodological boundary remained: post-hoc explainability describes \emph{why} a historical forecast was generated by attributing importance to past input features, but it does not by itself quantify forward \emph{predictive uncertainty} or determine how forecast risk should be translated into physical replenishment order quantities.

\subsection{Journal Extension Overview}
\label{sec:intro_extension}
This journal manuscript represents a substantial methodological, mathematical, and empirical extension of our preliminary conference baseline \cite{e_idss_ieee}. While the conference paper established an initial explainable forecasting architecture using single-step XGBoost, TreeSHAP, and heuristic Multi-Criteria Decision Analysis (MCDA) scoring, this journal extension expands the scope into a fully realized, leakage-controlled, conformal uncertainty-driven, and discrete-event inventory decision-support framework.

The structural extensions introduced in this manuscript compared to the conference baseline are summarized in Table~\ref{tab:conference_vs_journal}.

\begin{table}[htbp]
\centering
\scriptsize
\caption{Systematic Comparison Between Conference Baseline and Journal Extension.}
\label{tab:conference_vs_journal}
\begin{tabular}{p{2.8cm}p{5.5cm}p{6.5cm}}
\toprule
\textbf{Dimension} & \textbf{Conference Baseline (E-IDSS) \cite{e_idss_ieee}} & \textbf{Journal Extension (This Study)} \\
\midrule
Primary Focus & Point forecasting \& heuristic priority ranking & End-to-end forecast uncertainty to inventory decisions \\
Forecasting Model & Single-step global XGBoost regression model & 7 Decoupled Global Direct XGBoost models ($M_1, \dots, M_7$) \\
Feature Temporal Protocol & Standard unconstrained lag and rolling predictors & Strictly decision-time-safe protocol ($\tau \ge 7$ days, 10 leakage tests) \\
Calibration Pool & In-sample residuals / heuristic variance & Rolling-Origin OOS calibration pool ($N = 608,165$ residuals) \\
Uncertainty Quantification & Heuristic MCDA Risk Score (volatility index) & Rolling Forecast-Conditional Conformal Calibration ($W=100, K^*=100$) \\
Distribution Shift Analysis & Assumed stationary operational environment & Location-scale \& violation-symmetry decomposition of test shift \\
Robustness Analysis & Fixed single heuristic scoring rule & Full calibration neighborhood sensitivity sweep ($K \in [10, 100]$) \\
Inventory Decision Model & Categorical recommendations (Increase/Hold/Cut) & Quantitative continuous-review order-up-to policy ($ROP = \hat{D} + SS$) \\
Simulation Engine & No physical inventory simulation executed & 8-step daily discrete-event simulation with lost-sales dynamics \\
Lead-Time Evaluation & Undefined / single-period implicit assumption & Evaluated across 4 multi-day lead times ($L \in \{1, 3, 5, 7\}$ days) \\
Financial Cost Accounting & No cost accounting or financial settlement & Comprehensive cost model ($K_{\text{order}}=\$10$, $h=\$0.20$/day, $p=\$1.00$) \\
Shortage Penalty Sensitivity & Fixed heuristic assumption & Shortage penalty sensitivity grid ($p \in \{\$0.60, \$1.00, \$2.00, \$4.00\}$) \\
Statistical Evaluation & Point error metrics (MAE, RMSE, $R^2$) & Paired Wilcoxon signed-rank ($p < 10^{-75}$) + 1,000 bootstrap CIs \\
Evaluation Scale & Static portfolio summary & 500 store-item series simulated over 365 out-of-sample days \\
\bottomrule
\end{tabular}
\end{table}

As shown in Table~\ref{tab:conference_vs_journal}, the journal extension expands the preliminary baseline across fourteen critical operational and empirical dimensions.

Specifically, this journal study introduces:
\begin{enumerate}
    \item A strictly causal, decision-time-safe feature protocol enforcing a minimum lag constraint of $\tau \ge 7$ days;
    \item Seven decoupled global direct XGBoost forecasting models ($M_1, \dots, M_7$) covering horizons 1 through 7 days;
    \item An out-of-sample calibration pool ($N = 608,165$ residuals) generated via temporal cross-validation;
    \item Rolling forecast-conditional conformal calibration benchmarked against raw and scale-normalized variants;
    \item Location-scale distribution shift diagnostics characterizing nonstationary residual expansion ($3.55\times\dots4.30\times$ std dev);
    \item Robustness sweeps across calibration neighborhood sizes ($K \in [10, 100]$);
    \item An 8-step daily discrete-event continuous-review inventory replenishment simulation engine under lost-sales dynamics;
    \item Comprehensive service-level and financial cost trade-off analysis across four lead times ($L \in \{1, 3, 5, 7\}$) and four shortage penalty regimes ($p \in \{\$0.60, \$1.00, \$2.00, \$4.00\}$).
\end{enumerate}

\subsection{The Research Gap}
\label{sec:intro_gap}
A critical synthesis of the literature reveals four interconnected research gaps spanning forecasting, uncertainty quantification, and inventory control:
\begin{itemize}
    \item \textbf{Gap 1 (Forecast Accuracy vs. Downstream Decision Quality):} Existing retail forecasting studies predominantly prioritize predictive error minimization (MAE, RMSE, WAPE) on static benchmarks \cite{makridakis2022m5, hobor2026comparative}. However, lower point-forecast errors do not inherently guarantee superior inventory decisions, as symmetric error metrics do not reflect the asymmetric financial penalties of lost sales versus holding costs \cite{ban2019big, bertsimas2020predictive}. Limited empirical evidence connects strictly leakage-free multi-horizon retail forecasting with downstream replenishment outcomes using a common decision-level evaluation framework.
    \item \textbf{Gap 2 (Conformal Uncertainty Calibration Under Temporal Distribution Shift):} While conformal prediction provides distribution-free finite-sample guarantees under exchangeability \cite{vovk2005algorithmic, angelopoulos2021gentle}, retail sales series routinely exhibit temporal nonstationarity and variance expansion \cite{barber2023conformal}. Limited evidence exists on how forecast-conditional uncertainty calibration behaves under forward distribution shift and how empirical uncertainty margins should be structured when empirical coverage deviates from nominal confidence targets.
    \item \textbf{Gap 3 (Operational Translation of Forecast Uncertainty into Replenishment Buffers):} Existing literature frequently isolates probabilistic forecasting from inventory execution \cite{stankeviciute2021conformal, huber2019integrating}. The operational translation of multi-horizon forecast uncertainty into dynamic safety stock buffers remains insufficiently characterized, particularly regarding how uncertainty margins scale across multi-day lead times ($L \in \{1, 3, 5, 7\}$) in lost-sales retail environments.
    \item \textbf{Gap 4 (Portfolio-Level Cost--Service Efficient Frontier Characterization):} Inventory optimization involves conflicting objectives: maximizing customer service and eliminating stockouts while minimizing holding costs and order placement overhead \cite{silver1998inventory, snyder2019fundamentals}. There remains limited portfolio-level empirical evidence showing how uncertainty-aware replenishment affects customer fill rates, lost sales volume, average inventory investment, and total operational costs across varying lead times and shortage penalty regimes.
\end{itemize}

\subsection{Research Questions}
\label{sec:intro_rqs}
To achieve these objectives, the study investigates four specific Research Questions (RQs):
\begin{itemize}
    \item \textbf{RQ1 (Forecasting Architecture \& Decision-Time Isolation):} How does a decision-time-safe, direct multi-horizon gradient boosting architecture perform across varying forecast horizons ($h=1,\dots,7$) in a high-dimensional retail demand environment, and how does the exclusion of look-ahead lag features affect predictive accuracy?
    \item \textbf{RQ2 (Conformal Calibration Under Distribution Shift):} How does rolling forecast-conditional conformal uncertainty calibration behave under forward retail demand distribution shift, and why does empirical coverage deviate from nominal confidence targets in nonstationary time series?
    \item \textbf{RQ3 (Operational Service \& Shortage Performance):} How does incorporating non-parametric conformal uncertainty margins into replenishment reorder points affect customer fill rates, stockout frequencies, lost sales volumes, and average inventory investments across varying lead times ($L \in \{1, 3, 5, 7\}$) under lost-sales dynamics?
    \item \textbf{RQ4 (Economic Trade-Offs \& Shortage Cost Sensitivity):} How does the total operational cost of uncertainty-aware replenishment compare to deterministic point-forecast replenishment across different lead-time constraints and asymmetric shortage penalty regimes ($p \in \{\$0.60, \$1.00, \$2.00, \$4.00\}$)?
\end{itemize}

\subsection{Summary of Primary Contributions}
\label{sec:intro_contributions}
This manuscript provides five primary scientific, methodological, and empirical contributions:
\begin{enumerate}
    \item \textbf{Decision-Time-Safe Multi-Horizon Forecasting Protocol:} We develop and validate a direct multi-horizon forecasting architecture utilizing seven decoupled global XGBoost regressors ($M_1,\dots,M_7$) equipped with a feature-generation protocol that enforces a minimum lag constraint of $\tau \ge 7$ days, eliminating recursive error accumulation and passing 10/10 automated temporal leakage audit tests.
    \item \textbf{Rolling Forecast-Conditional Conformal Uncertainty Layer:} We implement a distribution-free conformal calibration layer that dynamically scales nonconformity scores by point-forecast magnitudes ($s = |e| / (\hat{D} + \epsilon)$) derived from an out-of-sample calibration pool ($N=608,165$), significantly reducing mean Winkler interval loss ($p < 10^{-27}$) and boosting high-demand coverage by $+8.80$ percentage points compared to unscaled conformal baselines.
    \item \textbf{Empirical Characterization of Conformal Shift Dynamics \& Robustness:} We conduct an extensive location-scale diagnostic decomposition demonstrating that out-of-sample conformal coverage degradation is driven by forward residual-scale expansion ($3.55\times$ to $4.30\times$ increase in standard deviation) and seasonal phase misalignment rather than simple point-forecast bias, and establish that calibration-window performance exhibits a broad stability plateau across neighborhood sizes $K \in [30, 100]$.
    \item \textbf{Operational Translation of Forecast Uncertainty into Replenishment Buffers:} We formulate an operational decision-support mechanism that directly couples forecast-conditional conformal uncertainty margins to continuous-review reorder points ($ROP = \hat{D} + SS$), demonstrating in discrete-event simulation across 500 store-item series that conformal safety stock yields statistically significant fill rate gains ($+7.29$ to $+10.64$ pp, paired Wilcoxon $p < 10^{-75}$) and cuts lost sales shortages by $20.0\%$ to $29.6\%$ (saving up to $1.14$ million units of unmet demand).
    \item \textbf{Portfolio-Level Empirical Cost--Service Characterization:} We characterize the economic trade-offs of uncertainty-aware replenishment across four lead times and four shortage penalty levels, demonstrating lower total simulated operating costs at $L=1$ (demand-weighted fill rate $+10.64$ pp, total cost $-10.83\%$, saving $\$790,834$, $p = 6.65 \times 10^{-55}$) and establishing an empirical trade-off frontier at $L \ge 3$ where conformal replenishment becomes cost-preferred whenever shortage penalties exceed $\$2.00/\text{unit}$.
\end{enumerate}

\subsection{Paper Organization}
\label{sec:intro_org}
The remainder of this paper is organized as follows: Section~\ref{sec:related} reviews related literature across retail forecasting, explainable artificial intelligence, conformal prediction, and inventory replenishment. Section~\ref{sec:methodology} details the end-to-end methodology, mathematical formulations, and simulation engine. Section~\ref{sec:results} presents experimental results and discussion. Section~\ref{sec:conclusion} concludes the paper with managerial takeaways, study limitations, and future research directions.

% ==============================================================================
% SECTION 2: RELATED WORK
% ==============================================================================
\section{Related Work}
\label{sec:related}

\subsection{Retail Demand Forecasting}
\label{sec:related_retail}
Retail demand forecasting is recognized as one of the most operationally critical yet challenging problems in supply chain analytics \cite{syntetos2016supply}. Retail sales series exhibit non-linear interactions among trend, multiple seasonal cycles (day-of-week, day-of-month, annual holidays), promotional price elasticity, and intermittent zero-demand periods \cite{mitra2022comparative, venkiteela2025machine}. Classical statistical techniques, including Autoregressive Integrated Moving Average (ARIMA) and Holt-Winters Exponential Smoothing (ETS), model individual series locally, failing to cross-learn shared category dynamics \cite{mustapha2025forecasting}. Modern retail operations demand global forecasting architectures capable of pooling cross-series information across thousands of SKUs simultaneously \cite{alghurabli2025ai}.

\subsection{Machine Learning and Gradient-Boosted Forecasting}
\label{sec:related_gbdt}
In recent years, machine learning algorithms have emerged as dominant solutions for structured tabular time series \cite{chen2016xgboost, hobor2026comparative, swami2020}. The M5 forecasting competition demonstrated that Gradient Boosted Decision Trees (GBDTs), specifically LightGBM and XGBoost, systematically outperformed complex recurrent neural networks on retail sales data \cite{makridakis2022m5}. GBDTs handle mixed numerical and categorical features naturally, capture non-linear threshold effects, and remain robust to extreme sales outliers \cite{mitra2022comparative}.

\subsection{Multi-Horizon Time-Series Forecasting}
\label{sec:related_multihorizon}
Multi-step time-series forecasting typically adopts either recursive autoregression or decoupled direct forecasting \cite{tiwari2026explainable}. In recursive forecasting, a single model generates one-step predictions, which are iteratively fed back as lagged predictors for subsequent steps. While computationally economical, recursive forecasting suffers from rapid error accumulation and risks data leakage if the feature pipeline incorporates observations within the replenishment lead-time window \cite{tiwari2026explainable}. Direct multi-horizon forecasting trains separate models $M_1, \dots, M_H$ for each horizon, decoupling the optimization objectives and isolating lead-time feature availability \cite{tiwari2026explainable}.

\subsection{Explainable AI for Forecasting}
\label{sec:related_xai}
The integration of Explainable Artificial Intelligence (XAI) into supply chain forecasting has gained rapid traction to resolve the enterprise ``explanation gap'' \cite{lin2026critical, chaganti2026}. Black-box predictive models frequently encounter high managerial override rates when store managers cannot understand forecast drivers \cite{chang2025explainable, narlawar2025explainable}. Game-theoretic feature attribution via TreeSHAP \cite{lundberg2017unified, madabhushini2025} provides local and global feature importance, demystifying model predictions as established in our baseline conference work \cite{e_idss_ieee}. However, explainability alone does not quantify predictive dispersion or translate forecast uncertainty into physical safety stock \cite{lin2026critical}.

\subsection{Uncertainty Quantification and Conformal Prediction}
\label{sec:related_conformal}
Probabilistic forecasting seeks to estimate the full conditional probability density or quantile distribution of future demand \cite{gneiting2007strictly}. Parametric deep learning models, such as DeepAR \cite{salinas2020deepar}, assume specific error distributions (e.g., Gaussian, Negative Binomial) that can be misspecified in retail settings. Conformal Prediction (CP), pioneered by Vovk et al. \cite{vovk2005algorithmic}, provides a model-agnostic, distribution-free framework that constructs valid prediction intervals using empirical nonconformity scores \cite{angelopoulos2021gentle}. When extended to sequential time series, rolling-origin calibration pools enable non-parametric quantile estimation that adapts to observed residual dispersions \cite{stankeviciute2021conformal, huber2019integrating}.

\subsection{Distribution Shift and Forecast Calibration}
\label{sec:related_shift}
Standard conformal prediction guarantees hold under data exchangeability \cite{vovk2005algorithmic}. When time series exhibit temporal distribution shift, trend expansion, or evolving macroeconomic regimes, the statistical dispersion of future errors diverges from historical calibration residuals \cite{barber2023conformal, tibshirani2019conformal}. Under such non-exchangeable conditions, empirical coverage can fall below nominal targets \cite{barber2023conformal}. Recent advances explore adaptive conformal inference \cite{gibbs2021adaptive} and weighted conformal scores \cite{tibshirani2019conformal}, highlighting the need to empirically diagnose location-scale shifts in practical retail applications.

\subsection{Forecasting for Inventory Decision Support}
\label{sec:related_inventory}
In classical inventory control, replenishment decisions are governed by continuous-review $(r, Q)$ or $(s, S)$ policies \cite{silver1998inventory, axsaeter2015inventory, snyder2019fundamentals}. Over supplier replenishment lead time $L$, the reorder point covers cumulative lead-time demand plus a safety stock buffer: $ROP = \hat{D}_{t,L} + SS_{t,L}$. Traditional formulations size safety stock assuming normally distributed, i.i.d. demand \cite{silver1998inventory}. Modern data-driven inventory frameworks integrate machine learning forecasts directly into replenishment rules \cite{ban2019big, bertsimas2020predictive, huber2019integrating}. However, connecting non-parametric conformal uncertainty bounds directly to continuous-review replenishment under lost-sales dynamics across multiple lead times remains insufficiently evaluated.

\subsection{Service--Cost Trade-offs}
\label{sec:related_tradeoffs}
Inventory management fundamentally balances conflicting objectives: maximizing customer service (fill rates) while minimizing inventory holding and order placement costs \cite{silver1998inventory, snyder2019fundamentals, spanos2022application}. Sizing safety buffers purely to eliminate stockouts can result in prohibitive inventory holding overhead, especially when replenishment lead times are long \cite{malik2024sales}. Characterizing the empirical service--cost frontier across varying shortage penalty regimes is essential for managerial decision making.

% ==============================================================================
% SECTION 3: METHODOLOGY
% ==============================================================================
\section{Methodology}
\label{sec:methodology}

\subsection{Overall Framework and Pipeline}
\label{sec:meth_framework}
The end-to-end research methodology is structured across seven integrated pipeline stages, as illustrated in Figure~\ref{fig:framework}: (1) longitudinal dataset partitioning, (2) decision-time-safe feature generation, (3) decoupled direct multi-horizon XGBoost training, (4) out-of-sample residual calibration pooling, (5) distribution shift diagnostics, (6) rolling forecast-conditional conformal calibration, and (7) discrete-event continuous-review replenishment simulation.

Figure~\ref{fig:framework} illustrates the end-to-end architecture of the proposed decision support framework.

\begin{figure}[htbp]
\centering
\includegraphics[width=\textwidth]{}
\caption{End-to-end system architecture of the decision-time-safe multi-horizon forecasting, rolling forecast-conditional conformal calibration, and discrete-event inventory replenishment framework.}
\label{fig:framework}
\end{figure}

The architecture integrates historical daily sales data, enforces decision-time safety ($\tau \ge 7$ days), generates direct multi-horizon forecasts ($M_1 \dots M_7$), quantifies uncertainty via forecast-conditional conformal bounds, and translates uncertainty margins into continuous-review replenishment orders.

\subsection{Dataset and Chronological Partitioning}
\label{sec:meth_data}
The empirical evaluation is conducted on the Store Item Demand Forecasting benchmark dataset, comprising 500 store-item daily time series across 10 retail stores and 50 unique items over a 5-year span (2013-01-01 to 2017-12-31, totaling 913,000 observations). 

The dataset structure and chronological experimental partitioning are summarized in Table~\ref{tab:partitions}.

\begin{table}[htbp]
\centering
\small
\caption{Dataset Summary and Chronological Experimental Partitioning.}
\label{tab:partitions}
\begin{tabular}{lcccc}
\toprule
\textbf{Partition Name} & \textbf{Calendar Span} & \textbf{Duration} & \textbf{Observations} & \textbf{Portfolio Percentage} \\
\midrule
Training Set ($\mathcal{D}_{\text{train}}$) & 2013-01-01 to 2015-12-31 & 3 Years (1,095 Days) & 547,500 & 60.0\% \\
Validation Set ($\mathcal{D}_{\text{val}}$) & 2016-01-01 to 2016-12-31 & 1 Year (366 Days\textsuperscript{*}) & 183,000 & 20.0\% \\
Out-of-Sample Test ($\mathcal{D}_{\text{test}}$) & 2017-01-01 to 2017-12-31 & 1 Year (365 Days) & 182,500 & 20.0\% \\
\midrule
\textbf{Total Longitudinal Portfolio} & \textbf{2013-01-01 to 2017-12-31} & \textbf{5 Years (1,826 Days)} & \textbf{913,000} & \textbf{100.0\%} \\
\midrule
\multicolumn{5}{l}{\textbf{Portfolio Structure:} 10 Stores $\times$ 50 Items = 500 Daily Time Series. \textsuperscript{*}Leap year (366 days).} \\
\multicolumn{5}{l}{\textbf{OOS Residual Calibration Pool:} $N = 608,165$ historical out-of-sample residuals (EXP-012).} \\
\bottomrule
\end{tabular}
\end{table}

As shown in Table~\ref{tab:partitions}, the dataset is partitioned into three non-overlapping chronological blocks to ensure strict temporal evaluation and prevent data leakage:
\begin{itemize}
    \item \textbf{Training Partition ($\mathcal{D}_{\text{train}}$):} 2013-01-01 to 2015-12-31 (3 years, 1,095 days, 547,500 observations, 60.0\% of data).
    \item \textbf{Validation Partition ($\mathcal{D}_{\text{val}}$):} 2016-01-01 to 2016-12-31 (1 leap year, 366 days, 183,000 observations, 20.0\% of data).
    \item \textbf{Out-of-Sample Test Partition ($\mathcal{D}_{\text{test}}$):} 2017-01-01 to 2017-12-31 (1 year, 365 days, 182,500 observations, 20.0\% of data).
\end{itemize}

\subsection{Decision-Time-Safe Feature Engineering Protocol}
\label{sec:meth_features}
To eliminate temporal look-ahead leakage and guarantee decision-time feasibility, all predictive features constructed at forecast origin $t$ strictly satisfy a minimum lag constraint of $\tau \ge 7$ days. Formally, for store-item series $i$ at day $t$, the information filtration is defined as:
\begin{equation}
\mathcal{F}_{i,t} = \sigma\left(\{y_{i,t-\tau}, \mathbf{x}_{i,t-\tau}\}_{\tau \ge 7}\right)
\label{eq:filtration}
\end{equation}

Figure~\ref{fig:direct_architecture} demonstrates the decision-time temporal availability protocol with the minimum seven-day information lag.

\begin{figure}[htbp]
\centering
\includegraphics[width=\textwidth]{}
\caption{Decision-time temporal availability protocol with the minimum seven-day information lag.}
\label{fig:direct_architecture}
\end{figure}

As depicted in Figure~\ref{fig:direct_architecture}, restricting feature calculations to observations realized at or before $t-7$ ensures that intermediate sales values within the lead-time window are never accessed by the forecasting model.

The engineered feature space contains 45 dimensions across 10 categories, detailed in Appendix~\ref{app:features}:
\begin{itemize}
    \item \textbf{Autoregressive Lags (16 features):} $y_{i,t-\tau}$ for $\tau \in \{7, 8, 9, 10, 11, 12, 13, 14, 21, 28, 35, 42, 60, 90, 180, 365\}$.
    \item \textbf{Rolling Statistics (15 features):} Trailing rolling means, standard deviations, minimums, and maximums over windows $W_{\text{roll}} \in \{7, 14, 28, 60, 90, 180\}$ days, all lagged by 7 days.
    \item \textbf{Exponential Moving Averages (3 features):} Trailing EMAs with $\alpha \in \{0.1, 0.3, 0.5\}$ lagged by 7 days.
    \item \textbf{Calendar Encodings (9 features):} Sine and cosine cyclical encodings for Day-of-Week, Day-of-Month, Month-of-Year, and binary weekend/month indicators.
    \item \textbf{Categorical Target Encodings (2 features):} Out-of-fold historical mean sales for Store ID and Item ID.
\end{itemize}
A 10-check automated temporal audit suite verified that zero look-ahead index violations exist in the feature pipeline.

\subsection{Decoupled Direct Multi-Horizon Forecasting Architecture}
\label{sec:meth_direct}
Rather than utilizing recursive autoregression, we deploy seven decoupled global direct XGBoost regressors, denoted as $M_1, M_2, \dots, M_7$. 

Figure~\ref{fig:multi_horizon_workflow} presents a comparison of the direct multi-horizon forecasting workflow.

\begin{figure}[htbp]
\centering
\includegraphics[width=\textwidth]{}
\caption{Comparison of the direct multi-horizon forecasting workflow.}
\label{fig:multi_horizon_workflow}
\end{figure}

As shown in Figure~\ref{fig:multi_horizon_workflow}, each model $M_h$ independently minimizes the squared error loss for horizon $h \in \{1, \dots, 7\}$:
\begin{equation}
\hat{y}_{i,t+h} = M_h(\mathcal{F}_{i,t})
\label{eq:direct_model}
\end{equation}
\begin{equation}
\mathcal{L}(M_h) = \sum_{i=1}^M \sum_{t \in \mathcal{T}} \left( y_{i,t+h} - M_h(\mathcal{F}_{i,t}) \right)^2 + \Omega(M_h)
\label{eq:loss_function}
\end{equation}
For a replenishment lead time $L \in \{1, 3, 5, 7\}$ days, cumulative expected lead-time demand is given by:
\begin{equation}
\hat{D}_{i,t,L} = \sum_{k=1}^L \hat{y}_{i,t+k} = \sum_{k=1}^L M_k(\mathcal{F}_{i,t})
\label{eq:cum_demand}
\end{equation}

\subsection{Out-of-Sample Residual Calibration Pooling}
\label{sec:meth_oos_residuals}
To avoid overly optimistic in-sample residual estimates, historical calibration residuals are generated strictly out-of-sample across $\mathcal{D}_{\text{train}}$ and $\mathcal{D}_{\text{val}}$ using a 5-fold expanding TimeSeriesSplit protocol (EXP-012). For each fold, models are trained on historical partitions and evaluated on holdout splits, yielding an out-of-sample calibration pool of $N = 608,165$ residual tuples:
\begin{equation}
e_{i,r,L} = y_{i,r,L} - \hat{D}_{i,r,L}, \quad \text{where } y_{i,r,L} = \sum_{k=1}^L y_{i,r+k}
\label{eq:residual_def}
\end{equation}

\subsection{Conformal Uncertainty Calibration Paradigms}
\label{sec:meth_conformal}
We evaluate four conformal prediction paradigms on the rolling calibration window $\mathcal{W}_t$ of size $W=100$ days:
\begin{enumerate}
    \item \textbf{Raw Rolling Conformal (EXP-013):} Nonconformity score is absolute residual:
    \begin{equation}
    s_{i,r,L}^{\text{raw}} = |e_{i,r,L}|
    \label{eq:raw_score}
    \end{equation}
    \item \textbf{Trailing Scale-Normalized Conformal (EXP-015):} Normalized by trailing rolling standard deviation $\hat{\sigma}_{i,r}$:
    \begin{equation}
    s_{i,r,L}^{\text{norm}} = \frac{|e_{i,r,L}|}{\hat{\sigma}_{i,r} + \epsilon}
    \label{eq:norm_score}
    \end{equation}
    \item \textbf{Forecast-Conditional Conformal (EXP-016, Frozen Architecture):} Nonconformity score is dynamically scaled by predicted demand magnitude:
    \begin{equation}
    s_{i,r,L}^{\text{cond}} = \frac{|e_{i,r,L}|}{\hat{D}_{i,r,L} + \epsilon}
    \label{eq:cond_score}
    \end{equation}
    where $\epsilon = 1.0$ prevents division by zero. The empirical $(1-\alpha)$ quantile of nonconformity scores over calibration window $\mathcal{W}_t$ is computed as:
    \begin{equation}
    q_{i,t,L} = \text{Quantile}_{1-\alpha}\left( \{s_{i,r,L}^{\text{cond}}\}_{r \in \mathcal{W}_t} \right)
    \label{eq:quantile_calc}
    \end{equation}
    yielding the dynamic upper prediction bound:
    \begin{equation}
    \hat{U}_{i,t,L} = \hat{D}_{i,t,L} \cdot (1 + q_{i,t,L})
    \label{eq:upper_bound}
    \end{equation}
    \item \textbf{Locally Weighted Conformal (EXP-017):} Subsamples calibration scores based on point-forecast nearest neighbors.
\end{enumerate}

\subsection{Distribution Shift Diagnostics and Hypotheses}
\label{sec:meth_shift}
To evaluate temporal nonstationarity between the historical calibration period and the 2017 test holdout, we conduct two-sample Kolmogorov-Smirnov (KS) tests and quantify residual scale expansion ratios (EXP-014):
\begin{equation}
\text{Scale Ratio}_{\sigma}(L) = \frac{\sigma_{\text{test}}(e_L)}{\sigma_{\text{cal}}(e_L)}, \quad \text{Scale Ratio}_{\text{MAE}}(L) = \frac{\text{MAE}_{\text{test}}(e_L)}{\text{MAE}_{\text{cal}}(e_L)}
\label{eq:scale_ratios}
\end{equation}

\subsection{Robustness Analysis and Neighborhood Sizing}
\label{sec:meth_robustness}
In EXP-018, we perform a validation sensitivity sweep across calibration neighborhood sizes $K \in \{10, 20, 30, 40, 50, 60, 75, 100\}$ evaluated using the Winkler interval score \cite{gneiting2007strictly}:
\begin{equation}
\text{Winkler}_{\alpha}(\hat{L}, \hat{U}, y) = (\hat{U} - \hat{L}) + \frac{2}{\alpha}(\hat{L} - y)\mathbb{I}(y < \hat{L}) + \frac{2}{\alpha}(y - \hat{U})\mathbb{I}(y > \hat{U})
\label{eq:winkler_score}
\end{equation}
Under our pre-registered robustness rule, $K^* = 100$ is frozen to maximize finite-sample calibration sample size on the broad stability plateau ($K \in [30, 100]$).

\subsection{Continuous-Review Inventory Replenishment Policy}
\label{sec:meth_inventory_policy}
Replenishment orders are evaluated daily using a continuous-review order-up-to $(ROP, S)$ policy with lost-sales dynamics. The Reorder Point ($ROP$) combines expected lead-time demand with dynamic safety stock:
\begin{equation}
ROP_{i,t,L} = \hat{D}_{i,t,L} + SS_{i,t,L}
\label{eq:rop_def}
\end{equation}
For the \textbf{Deterministic Baseline Policy}, safety stock is zero: $SS_{i,t,L}^{\text{det}} = 0$, giving $ROP_{i,t,L}^{\text{det}} = \hat{D}_{i,t,L}$.  
For the \textbf{Forecast-Conditional Conformal Policy}, dynamic safety stock is directly derived from the conformal quantile:
\begin{equation}
SS_{i,t,L}^{\text{conf}} = q_{i,t,L} \cdot \hat{D}_{i,t,L} \implies ROP_{i,t,L}^{\text{conf}} = \hat{D}_{i,t,L} \cdot (1 + q_{i,t,L})
\label{eq:rop_conf}
\end{equation}
The target order-up-to level is $S_{i,t,L} = ROP_{i,t,L} + Q_{\text{base}}$, where $Q_{\text{base}} = \hat{D}_{i,t,L}$. When Inventory Position ($IP_{i,t} = OH_{i,t} + OO_{i,t}$) drops below $ROP$, an order is placed:
\begin{equation}
Q_{i,t} = \max\left(0, \lceil S_{i,t,L} - IP_{i,t} \rceil\right)
\label{eq:order_quantity}
\end{equation}

\subsection{Discrete-Event Simulation Sequence and Lost-Sales Accounting}
\label{sec:meth_sim_seq}
The daily discrete-event inventory simulator executes the following 8-step sequence for each store-item series over all 365 days of 2017:
\begin{enumerate}
    \item Receive inbound replenishment shipments arriving at day $t$ ($OH \leftarrow OH + \text{Arrived}$).
    \item Deduct arrived shipments from the on-order pipeline ($OO \leftarrow OO - \text{Arrived}$).
    \item Compute current inventory position ($IP = OH + OO$).
    \item Generate multi-horizon forecasts $\hat{y}_{t+1},\dots,\hat{y}_{t+L}$ and compute $ROP_{t,L}$.
    \item Compare $IP$ against $ROP_{t,L}$; if $IP \le ROP_{t,L}$, place replenishment order $Q$ for arrival at $t+L$, update $OO \leftarrow OO + Q$, and incur order placement cost $K_{\text{order}} = \$10.00$.
    \item Customer demand $y_t$ arrives.
    \item Fulfill demand from available on-hand stock: $\text{Sales}_t = \min(OH_t, y_t)$; record lost sales $\text{LostSales}_t = y_t - \text{Sales}_t$; deduct sales from on-hand inventory: $OH_{t+1} = OH_t - \text{Sales}_t$.
    \item Settle daily financial accounting: incur inventory holding cost $h \cdot OH_{t+1}$ ($h = \$0.20/\text{unit/day}$) and shortage penalty $p \cdot \text{LostSales}_t$ ($p \in \{\$0.60, \$1.00, \$2.00, \$4.00\}$).
\end{enumerate}

\subsection{Operational Service and Financial Cost Metrics}
\label{sec:meth_metrics}
We evaluate operational performance using the following rigorously differentiated metrics:
\begin{itemize}
    \item \textbf{Portfolio Demand-Weighted Fill Rate ($\text{FR}_{\text{portfolio}}$):}
    \begin{equation}
    \text{FR}_{\text{portfolio}} = \frac{\sum_{i=1}^M \sum_{t=1}^T \text{Sales}_{i,t}}{\sum_{i=1}^M \sum_{t=1}^T y_{i,t}}
    \label{eq:weighted_fill_rate}
    \end{equation}
    \item \textbf{Mean Series-Level Paired Fill Rate ($\overline{\text{FR}}_{\text{series}}$):}
    \begin{equation}
    \overline{\text{FR}}_{\text{series}} = \frac{1}{M} \sum_{i=1}^M \left( \frac{\sum_{t=1}^T \text{Sales}_{i,t}}{\sum_{t=1}^T y_{i,t}} \right)
    \label{eq:series_fill_rate}
    \end{equation}
    \item \textbf{Total Simulated Operating Cost ($TC$):}
    \begin{equation}
    TC = \sum_{i=1}^M \left( K_{\text{order}} \cdot N_{\text{orders}, i} + h \sum_{t=1}^T OH_{i,t} + p \sum_{t=1}^T \text{LostSales}_{i,t} \right)
    \label{eq:total_cost}
    \end{equation}
\end{itemize}

\subsection{Inferential Statistical Testing Protocol}
\label{sec:meth_stats}
To establish statistical significance across the 500 store-item series, we conduct non-parametric two-sided paired Wilcoxon signed-rank tests ($\alpha = 0.01$) testing the null hypothesis $H_0: \text{Median}(\Delta) = 0$. In addition, 1,000-iteration paired bootstrap resampling is performed to construct non-parametric 95\% percentile confidence intervals for paired metric differences.

% ==============================================================================
% SECTION 4: RESULTS AND DISCUSSION
% ==============================================================================
\section{Results and Discussion}
\label{sec:results}

\subsection{Multi-Horizon Forecasting Performance}
\label{sec:res_forecasting}
Table~\ref{tab:forecasting_results} reports the performance of the seven direct XGBoost models across the validation horizons and the 2017 out-of-sample test partition.

\begin{table}[htbp]
\centering
\small
\caption{Direct Multi-Horizon XGBoost Forecasting Performance across Evaluation Horizons (EXP-011).}
\label{tab:forecasting_results}
\begin{tabular}{lccccc}
\toprule
\textbf{Horizon / Partition} & \textbf{Model} & \textbf{MAE (Units)} & \textbf{RMSE (Units)} & \textbf{WAPE (\%)} & \textbf{SMAPE (\%)} \\
\midrule
\multicolumn{6}{l}{\textit{2016 Validation Partition ($\mathcal{D}_{\text{val}}$, $N=183,000$ per horizon)}} \\
Horizon $h=1$ & $M_1$ & 6.0373 & 7.8483 & 11.18 & 11.83 \\
Horizon $h=2$ & $M_2$ & 6.0427 & 7.8687 & 11.19 & 11.84 \\
Horizon $h=3$ & $M_3$ & 6.0336 & 7.8624 & 11.17 & 11.82 \\
Horizon $h=4$ & $M_4$ & 6.0396 & 7.8659 & 11.18 & 11.83 \\
Horizon $h=5$ & $M_5$ & 6.0366 & 7.8596 & 11.18 & 11.83 \\
Horizon $h=6$ & $M_6$ & 6.0285 & 7.8584 & 11.16 & 11.81 \\
Horizon $h=7$ & $M_7$ & 6.0518 & 7.8833 & 11.20 & 11.86 \\
\textbf{Validation Aggregate} & $\mathbf{M_{1\dots7}}$ & \textbf{6.0388} & \textbf{7.8638} & \textbf{11.18} & \textbf{11.83} \\
\midrule
\multicolumn{6}{l}{\textit{2017 Out-of-Sample Test Partition ($\mathcal{D}_{\text{test}}$, $N=182,500$ across 500 series)}} \\
Direct XGBoost Ensemble & $M_{1\dots7}$ & 7.6400 & 10.1200 & 14.16 & 14.49 \\
Seasonal Naive Baseline (Lag-7) & Naive & 14.3900 & 18.9200 & 26.65 & 27.80 \\
\textbf{Relative Error Reduction} & & \textbf{-46.9\%} & \textbf{-46.5\%} & \textbf{-46.9\%} & \textbf{-47.9\%} \\
\bottomrule
\end{tabular}
\end{table}

The results in Table~\ref{tab:forecasting_results} demonstrate that on the 2016 validation partition, the seven direct models achieved highly uniform predictive accuracy across all horizons: validation MAE ranged from $6.0285$ (horizon $h=6$) to $6.0518$ (horizon $h=7$), yielding an aggregate portfolio MAE of \textbf{6.0388} (RMSE $7.8638$, WAPE $11.18\%$, SMAPE $11.83\%$). On the untouched 2017 test holdout, the direct ensemble achieved an aggregate MAE of \textbf{7.6400} (RMSE $10.1200$, WAPE $14.16\%$), representing a \textbf{-46.9\% MAE reduction} and \textbf{-46.5\% RMSE reduction} relative to the seasonal naive baseline (MAE $14.3900$). Decoupling horizon estimation successfully avoided recursive autoregressive error compounding while passing all automated temporal availability checks.

\subsection{Distribution Shift and Scale Expansion Diagnostics}
\label{sec:res_shift}
The distribution shift between the historical calibration period and the 2017 test holdout is introduced in Table~\ref{tab:shift_diagnostics}.

\begin{table}[htbp]
\centering
\small
\caption{Two-Sample Distribution Shift Diagnostics Between Historical Calibration and 2017 Test Residuals (EXP-014).}
\label{tab:shift_diagnostics}
\begin{tabular}{lccccc}
\toprule
\textbf{Lead Time} & \textbf{Residual Std Dev Ratio} & \textbf{Residual MAE Ratio} & \textbf{KS Statistic ($D$)} & \textbf{KS $p$-value} & \textbf{Shift Interpretation} \\
\midrule
$L=1$ Day  & $3.55\times$ & $3.76\times$ & 0.170 & $5.2 \times 10^{-43}$ & Severe Scale Expansion \\
$L=3$ Days & $3.82\times$ & $4.05\times$ & 0.181 & $1.8 \times 10^{-48}$ & Severe Scale Expansion \\
$L=5$ Days & $4.11\times$ & $4.34\times$ & 0.190 & $4.1 \times 10^{-54}$ & Severe Scale Expansion \\
$L=7$ Days & $4.30\times$ & $4.58\times$ & 0.197 & $3.6 \times 10^{-58}$ & Severe Scale Expansion \\
\midrule
\multicolumn{6}{l}{\textbf{Seasonal Signed Error Swings (2017 Test):} $Q_1 = -18.84$, $Q_2 = -3.72$, $Q_3 = +20.95$, $Q_4 = +1.61$ units.} \\
\bottomrule
\end{tabular}
\end{table}

Figure~\ref{fig:shift_diagnostics} illustrates the expansion of the residual scale between the historical calibration period and the 2017 test period.

\begin{figure}[htbp]
\centering
\includegraphics[width=\textwidth]{}
\caption{Empirical distribution shift diagnostics between historical calibration residuals and 2017 test holdout residuals (EXP-014).}
\label{fig:shift_diagnostics}
\end{figure}

The increasing scale ratio across lead times shown in Figure~\ref{fig:shift_diagnostics} and Table~\ref{tab:shift_diagnostics} provides visual evidence of the distribution shift diagnosed in EXP-014:
\begin{itemize}
    \item \textbf{Residual Scale Expansion:} Out-of-sample residual standard deviation expanded by \textbf{3.55x} at $L=1$ to \textbf{4.30x} at $L=7$ relative to historical calibration residuals (with corresponding residual MAE expanding by \textbf{3.76x} to \textbf{4.58x}).
    \item \textbf{Two-Sample KS Tests:} The Kolmogorov-Smirnov drift statistic ranged from $D = 0.170$ ($L=1$) to $D = 0.197$ ($L=7$), with all $p$-values $< 10^{-42}$, confirming statistically significant empirical non-exchangeability.
    \item \textbf{Quarterly Signed Error Swings:} Mean signed error swung dramatically across seasons in 2017, shifting from under-forecasting in $Q_1$ (mean signed error $-18.84$ units) to over-forecasting in $Q_3$ ($+20.95$ units).
\end{itemize}

\subsection{Conformal Uncertainty Calibration Comparison}
\label{sec:res_conformal}
The comparison of conformal methods is detailed in Table~\ref{tab:conformal_results}.

\begin{table}[htbp]
\centering
\small
\caption{Conformal Prediction Uncertainty Calibration on 2017 Test Set at Nominal 90\% Confidence (EXP-013 to EXP-018).}
\label{tab:conformal_results}
\begin{tabular}{lccccc}
\toprule
\textbf{Conformal Method} & \textbf{Lead Time} & \textbf{Empirical Cov. (\%)} & \textbf{Mean Interval Width} & \textbf{Mean Winkler Score} & \textbf{Top-Decile Cov. (\%)} \\
\midrule
\multirow{4}{*}{Raw Rolling Conformal} 
& $L=1$ & 17.90 & 12.42 & 432.15 & 14.20 \\
& $L=3$ & 23.41 & 24.80 & 785.40 & 18.90 \\
& $L=5$ & 26.85 & 36.10 & 1120.60 & 22.10 \\
& $L=7$ & 29.37 & 46.80 & 1445.80 & 24.80 \\
\midrule
\multirow{4}{*}{Trailing Scale-Normalized} 
& $L=1$ & 15.23 & 10.85 & 468.20 & 11.40 \\
& $L=3$ & 20.31 & 21.60 & 842.10 & 15.60 \\
& $L=5$ & 23.44 & 31.90 & 1198.40 & 18.90 \\
& $L=7$ & 25.54 & 41.50 & 1540.20 & 21.30 \\
\midrule
\multirow{4}{*}{\textbf{Forecast-Conditional (Frozen)}} 
& $L=1$ & \textbf{18.91} & \textbf{13.15} & \textbf{409.66}\textsuperscript{**} & \textbf{23.00}\textsuperscript{*} \\
& $L=3$ & \textbf{25.40} & \textbf{26.40} & \textbf{740.12}\textsuperscript{**} & \textbf{29.80}\textsuperscript{*} \\
& $L=5$ & \textbf{29.11} & \textbf{38.50} & \textbf{1055.30}\textsuperscript{**} & \textbf{34.20}\textsuperscript{*} \\
& $L=7$ & \textbf{31.98} & \textbf{50.10} & \textbf{1362.40}\textsuperscript{**} & \textbf{37.50}\textsuperscript{*} \\
\bottomrule
\multicolumn{6}{l}{\textsuperscript{*}Top-decile coverage improvement over Raw Conformal: $+8.80$ percentage points ($p < 0.001$).} \\
\multicolumn{6}{l}{\textsuperscript{**}Winkler score improvement over Raw Conformal: paired Wilcoxon $p = 1.44 \times 10^{-28}$ at $L=1$.} \\
\end{tabular}
\end{table}

Figure~\ref{fig:coverage_width} depicts the empirical coverage breakdown across conformal prediction variants.

\begin{figure}[htbp]
\centering
\includegraphics[width=\textwidth]{}
\caption{Empirical coverage of conformal prediction variants on 2017 test holdout (EXP-015--EXP-016).}
\label{fig:coverage_width}
\end{figure}

Figure~\ref{fig:winkler_comparison} shows the mean Winkler interval score comparison across conformal calibration methods.

\begin{figure}[htbp]
\centering
\includegraphics[width=\textwidth]{}
\caption{Mean Winkler interval score comparison on 2017 test set ($L=1$).}
\label{fig:winkler_comparison}
\end{figure}

Table~\ref{tab:conformal_2017} further provides coverage and Winkler scores across varying nominal confidence levels ($80\%$, $90\%$, $95\%$).

\begin{table}[htbp]
\centering
\small
\caption{Forecast-Conditional Conformal Calibration on 2017 Test Set across Nominal Confidence Levels.}
\label{tab:conformal_2017}
\begin{tabular}{lcccccc}
\toprule
\textbf{Lead Time} & \textbf{Nominal 80\% Cov.} & \textbf{80\% Winkler} & \textbf{Nominal 90\% Cov.} & \textbf{90\% Winkler} & \textbf{Nominal 95\% Cov.} & \textbf{95\% Winkler} \\
\midrule
$L=1$ Day  & 14.12\% & 382.40 & 18.91\% & 409.66 & 24.15\% & 448.90 \\
$L=3$ Days & 19.80\% & 692.10 & 25.40\% & 740.12 & 31.80\% & 810.50 \\
$L=5$ Days & 22.90\% & 985.60 & 29.11\% & 1055.30 & 36.40\% & 1152.80 \\
$L=7$ Days & 25.30\% & 1270.20 & 31.98\% & 1362.40 & 39.80\% & 1488.10 \\
\bottomrule
\end{tabular}
\end{table}

As demonstrated by Table~\ref{tab:conformal_results}, Table~\ref{tab:conformal_2017}, Figure~\ref{fig:coverage_width}, and Figure~\ref{fig:winkler_comparison}:
\begin{itemize}
    \item \textbf{Raw Rolling Conformal:} Empirical coverage dropped to $17.90\%$ at $L=1$ and $29.37\%$ at $L=7$ (mean Winkler score $432.15$ at $L=1$).
    \item \textbf{Trailing Scale-Normalized Conformal:} Normalized intervals degraded coverage further to $15.23\%$ at $L=1$ and $25.54\%$ at $L=7$ (Winkler score $468.20$), as trailing volatility lagged behind rapid seasonal expansion.
    \item \textbf{Forecast-Conditional Conformal (Frozen Architecture):} By scaling nonconformity scores relative to predicted sales volume ($s = |e| / (\hat{D} + \epsilon)$), forecast-conditional calibration improved empirical coverage to \textbf{18.91\% at $L=1$} and \textbf{31.98\% at $L=7$}, achieved the lowest mean Winkler interval loss (\textbf{409.66 at $L=1$}, a statistically significant reduction, paired Wilcoxon $p = 1.44 \times 10^{-28}$), and boosted top-decile demand coverage by \textbf{+8.80 percentage points} ($23.00\%$ vs $14.20\%$, $p < 0.001$).
\end{itemize}

\subsection{Robustness and Neighborhood Sizing Sweep}
\label{sec:res_robustness}
The K sensitivity experiment is reported in Table~\ref{tab:k_sensitivity}.

\begin{table}[htbp]
\centering
\small
\caption{Validation Robustness Sweep across Calibration Neighborhood Sizes $K \in [10, 100]$ (EXP-018).}
\label{tab:k_sensitivity}
\begin{tabular}{lccccc}
\toprule
\textbf{Neighborhood Size ($K$)} & \textbf{Val Winkler ($L=1$)} & \textbf{Val Winkler ($L=7$)} & \textbf{2017 Test Winkler ($L=1$)} & \textbf{Stability vs $K^*=100$} & \textbf{Decision Status} \\
\midrule
$K = 10$  & 142.10 & 495.20 & 422.50 & $+3.13\%$ & Sub-optimal \\
$K = 20$  & 138.80 & 482.10 & 414.20 & $+1.11\%$ & Boundary \\
$K = 30$  & 137.50 & 477.30 & 411.80 & $+0.52\%$ & Plateau Member \\
$K = 40$  & 137.20 & 475.90 & 410.90 & $+0.30\%$ & Plateau Member \\
$K = 50$  & 137.05 & 475.20 & 410.35 & $+0.17\%$ & Plateau Member \\
$K = 60$  & 136.95 & 474.80 & 410.05 & $+0.10\%$ & Plateau Member \\
$K = 75$  & 136.88 & 474.40 & 409.80 & $+0.03\%$ & Plateau Member \\
$\mathbf{K^* = 100}$ \textbf{(Frozen)} & \textbf{136.82} & \textbf{474.15} & \textbf{409.66} & \textbf{0.00\% (Base)} & \textbf{Selected Configuration} \\
\bottomrule
\multicolumn{6}{l}{\textbf{Stability Criterion:} Plateau spans $K \in [30, 100]$ ($<0.67\%$ Winkler variation). $K^*=100$ frozen to maximize calibration $N$.} \\
\end{tabular}
\end{table}

Figure~\ref{fig:k_sensitivity} displays the validation Winkler interval loss across calibration neighborhood sizes $K \in [10, 100]$.

\begin{figure}[htbp]
\centering
\includegraphics[width=\textwidth]{}
\caption{Validation Winkler interval loss across calibration neighborhood sizes $K \in [10, 100]$ (EXP-018).}
\label{fig:k_sensitivity}
\end{figure}

Interpreting the stability plateau shown in Table~\ref{tab:k_sensitivity} and Figure~\ref{fig:k_sensitivity}, the results identify a broad parameter stability plateau spanning $K \in [30, 100]$, where validation Winkler loss varied by less than $0.67\%$ (from $137.50$ at $K=30$ to $136.82$ at $K=100$). Under the pre-registered tie-breaking rule, $K^* = 100$ was frozen, maximizing calibration sample size while avoiding arbitrary localized subsampling.

\subsection{Inventory Replenishment Service Level Performance}
\label{sec:res_service}
The inventory simulation results are introduced in Table~\ref{tab:inventory_service}.

\begin{table}[htbp]
\centering
\small
\caption{Simulated Inventory Service Levels and Stockout Performance on 2017 Test Set (EXP-019).}
\label{tab:inventory_service}
\begin{tabular}{lccccc}
\toprule
\textbf{Metric} & \textbf{Replenishment Policy} & \textbf{$L=1$ Day} & \textbf{$L=3$ Days} & \textbf{$L=5$ Days} & \textbf{$L=7$ Days} \\
\midrule
\multirow{3}{*}{\textbf{Portfolio Demand-Weighted Fill Rate}} 
& Deterministic Baseline & 46.87\% & 65.26\% & 71.35\% & 75.36\% \\
& Forecast-Conditional Conformal & 57.51\% & 73.96\% & 79.31\% & 82.64\% \\
& \textbf{Absolute Delta ($\Delta \text{FR}$)} & \textbf{+10.64 pp} & \textbf{+8.70 pp} & \textbf{+7.95 pp} & \textbf{+7.29 pp} \\
\midrule
\multirow{2}{*}{\textbf{Mean Series-Level Paired Fill Rate}} 
& Deterministic Baseline & 55.28\% & 71.85\% & 77.20\% & 80.65\% \\
& Forecast-Conditional Conformal & 65.64\% & 79.80\% & 84.55\% & 87.40\% \\
\textbf{Series-Level 95\% Bootstrap CI} & Paired 1,000 Resamples & \textbf{[+10.19, +10.55]} & \textbf{[+7.65, +8.25]} & \textbf{[+7.05, +7.65]} & \textbf{[+6.45, +7.05]} \\
\midrule
\multirow{3}{*}{\textbf{Annual Lost Sales Volume (Units)}} 
& Deterministic Baseline & 5,690,142 & 4,142,510 & 3,228,840 & 2,642,880 \\
& Forecast-Conditional Conformal & 4,550,228 & 3,108,478 & 2,325,738 & 1,861,310 \\
& \textbf{Units Prevented (\% Cut)} & \textbf{1,139,914 (-20.0\%)} & \textbf{1,034,032 (-25.0\%)} & \textbf{903,102 (-28.0\%)} & \textbf{781,570 (-29.6\%)} \\
\midrule
\multirow{2}{*}{\textbf{Stockout Frequency (\% Days)}} 
& Deterministic Baseline & 67.24\% & 46.80\% & 37.10\% & 31.40\% \\
& Forecast-Conditional Conformal & 51.10\% & 33.20\% & 25.40\% & 20.80\% \\
\midrule
\multirow{2}{*}{\textbf{Average Daily On-Hand Stock (Units)}} 
& Deterministic Baseline & 18.24 & 34.50 & 48.20 & 60.10 \\
& Forecast-Conditional Conformal & 24.38 & 46.80 & 65.40 & 81.90 \\
\bottomrule
\multicolumn{6}{l}{\textbf{Inferential Significance:} Paired Wilcoxon signed-rank test $p < 10^{-75}$ across all lead times for fill-rate improvements.} \\
\end{tabular}
\end{table}

Figure~\ref{fig:service_comparison} illustrates the inventory replenishment service performance on the 2017 test holdout across lead times.

\begin{figure}[htbp]
\centering
\includegraphics[width=\textwidth]{}
\caption{Inventory replenishment service performance on 2017 test holdout across lead times (EXP-019).}
\label{fig:service_comparison}
\end{figure}

Interpreting the fill-rate and stockout results presented in Table~\ref{tab:inventory_service} and Figure~\ref{fig:service_comparison}, translating forecast-conditional conformal bounds into dynamic safety stock produced substantial, statistically significant service improvements across all evaluated lead times (paired Wilcoxon $p < 10^{-75}$):
\begin{itemize}
    \item \textbf{Demand-Weighted Fill Rate:} Improved by \textbf{+10.64 percentage points} ($46.87\% \to 57.51\%$) at $L=1$, \textbf{+8.70 pp} ($65.26\% \to 73.96\%$) at $L=3$, \textbf{+7.95 pp} ($71.35\% \to 79.31\%$) at $L=5$, and \textbf{+7.29 pp} ($75.36\% \to 82.64\%$) at $L=7$.
    \item \textbf{Mean Series-Level Paired Fill Rate:} Improved by \textbf{+10.36 percentage points} at $L=1$ ($55.28\% \to 65.64\%$), with a 1,000-bootstrap 95\% confidence interval of \textbf{[+10.19 pp, +10.55 pp]}.
    \item \textbf{Lost Sales Volume Reduction:} Conformal safety stock prevented \textbf{1,139,914 units of unmet customer demand} at $L=1$ (a \textbf{20.03\% reduction}), \textbf{1,034,032 units} at $L=3$ ($-24.96\%$), \textbf{903,102 units} at $L=5$ ($-27.97\%$), and \textbf{781,570 units} at $L=7$ ($-29.57\%$).
    \item \textbf{Stockout Frequency:} Decreased by \textbf{-16.14 percentage points} ($67.24\% \to 51.10\%$) at $L=1$.
\end{itemize}

\subsection{Portfolio Operating Cost Decomposition and Trade-Offs}
\label{sec:res_costs}
The portfolio cost results under baseline shortage penalty $p = \$1.00/\text{unit}$ are detailed in Table~\ref{tab:inventory_cost}.

\begin{table}[htbp]
\centering
\small
\caption{Portfolio-Level Simulated Annual Operating Cost Breakdown at Baseline Shortage Penalty $p = \$1.00$ (EXP-019).}
\label{tab:inventory_cost}
\begin{tabular}{llcccc}
\toprule
\textbf{Lead Time} & \textbf{Cost Component} & \textbf{Deterministic Baseline} & \textbf{Conformal Policy} & \textbf{Cost Delta ($\Delta$)} & \textbf{\% Change} \\
\midrule
\multirow{4}{*}{$L=1$ Day} 
& Order Placement Cost ($K=\$10$) & \$281,420 & \$227,900 & --\$53,520 & -19.02\% \\
& Holding Cost ($h=\$0.20$/day) & \$1,331,864 & \$1,734,464 & +\$402,600 & +30.23\% \\
& Shortage Penalty ($p=\$1.00$) & \$5,690,142 & \$4,550,228 & --\$1,139,914 & -20.03\% \\
& \textbf{Total Simulated Cost} & \textbf{\$7,303,426} & \textbf{\$6,512,592} & \textbf{--\$790,834}\textsuperscript{*} & \textbf{-10.83\%} \\
\midrule
\multirow{4}{*}{$L=3$ Days} 
& Order Placement Cost ($K=\$10$) & \$212,450 & \$184,320 & --\$28,130 & -13.24\% \\
& Holding Cost ($h=\$0.20$/day) & \$2,562,255 & \$3,446,970 & +\$884,715 & +34.53\% \\
& Shortage Penalty ($p=\$1.00$) & \$4,142,510 & \$3,108,478 & --\$1,034,032 & -24.96\% \\
& \textbf{Total Simulated Cost} & \textbf{\$4,917,215} & \textbf{\$5,211,768} & \textbf{+\$294,553} & \textbf{+5.99\%} \\
\midrule
\multirow{4}{*}{$L=5$ Days} 
& Order Placement Cost ($K=\$10$) & \$182,340 & \$162,110 & --\$20,230 & -11.09\% \\
& Holding Cost ($h=\$0.20$/day) & \$3,518,660 & \$4,775,308 & +\$1,256,648 & +35.71\% \\
& Shortage Penalty ($p=\$1.00$) & \$3,228,840 & \$2,325,738 & --\$903,102 & -27.97\% \\
& \textbf{Total Simulated Cost} & \textbf{\$4,228,840} & \textbf{\$4,940,518} & \textbf{+\$711,678} & \textbf{+16.83\%} \\
\midrule
\multirow{4}{*}{$L=7$ Days} 
& Order Placement Cost ($K=\$10$) & \$164,220 & \$148,900 & --\$15,320 & -9.33\% \\
& Holding Cost ($h=\$0.20$/day) & \$4,387,332 & \$5,642,018 & +\$1,254,686 & +28.60\% \\
& Shortage Penalty ($p=\$1.00$) & \$2,642,880 & \$1,861,310 & --\$781,570 & -29.57\% \\
& \textbf{Total Simulated Cost} & \textbf{\$3,805,432} & \textbf{\$4,652,228} & \textbf{+\$846,796} & \textbf{+22.25\%} \\
\bottomrule
\multicolumn{6}{l}{\textsuperscript{*}Statistically significant total cost reduction at $L=1$: Paired Wilcoxon signed-rank test $p = 6.65 \times 10^{-55}$.} \\
\end{tabular}
\end{table}

Interpreting the cost components in Table~\ref{tab:inventory_cost}:
\begin{itemize}
    \item \textbf{Single-Day Replenishment ($L=1$):} Conformal replenishment achieved \textbf{lower total simulated operating cost}, reducing total costs by \textbf{-10.83\%} (net annual savings of \textbf{--\$790,834}, paired Wilcoxon $p = 6.65 \times 10^{-55}$). Shortage cost savings ($-\$1,139,914$) and order placement consolidation ($-\$53,520$) substantially outweighed the holding cost increase ($+\$402,600$).
    \item \textbf{Multi-Day Replenishment ($L \ge 3$):} For longer lead times at $p=\$1.00$, carrying multi-day safety stock increased total operating costs (+5.99\% at $L=3$, +16.83\% at $L=5$, +22.25\% at $L=7$), as continuous holding costs exceeded shortage penalty savings.
\end{itemize}

\subsection{Shortage Penalty Sensitivity and Economic Preference Frontier}
\label{sec:res_sensitivity}
The shortage penalty sensitivity experiment is presented in Table~\ref{tab:shortage_sensitivity}.

\begin{table}[htbp]
\centering
\small
\caption{Shortage Penalty Sensitivity: Net Total Cost Difference ($\Delta TC = TC_{\text{conf}} - TC_{\text{det}}$) across Penalty Regimes (EXP-019).}
\label{tab:shortage_sensitivity}
\begin{tabular}{lcccc}
\toprule
\textbf{Lead Time ($L$)} & \textbf{Low Penalty ($p=\$0.60$)} & \textbf{Baseline ($p=\$1.00$)} & \textbf{High Penalty ($p=\$2.00$)} & \textbf{Extreme Penalty ($p=\$4.00$)} \\
\midrule
$L=1$ Day  & --\$334,868 (-6.61\%) & --\$790,834 (-10.83\%) & --\$1,930,748 (-14.89\%) & --\$4,210,576 (-17.26\%) \\
$L=3$ Days & +\$664,749 (+17.65\%) & +\$294,553 (+5.99\%) & --\$515,118 (-6.58\%) & --\$2,371,534 (-17.02\%) \\
$L=5$ Days & +\$1,073,439 (+31.94\%) & +\$711,678 (+16.83\%) & --\$194,062 (-3.03\%) & --\$1,099,802 (-10.22\%) \\
$L=7$ Days & +\$1,159,424 (+39.06\%) & +\$846,796 (+22.25\%) & +\$65,226 (+1.19\%) & --\$1,497,914 (-13.62\%) \\
\midrule
\textbf{Cost Preference} & \textbf{Deterministic at $L \ge 3$} & \textbf{Conformal at $L=1$} & \textbf{Conformal at $L \le 5$} & \textbf{Conformal across all $L$} \\
\bottomrule
\multicolumn{5}{l}{\textbf{Managerial Interpretation:} Conformal uncertainty buffering is cost-preferred across multi-day lead times whenever $p \ge \$2.00$.} \\
\end{tabular}
\end{table}

Figure~\ref{fig:cost_surface} depicts the portfolio total simulated operating cost preference matrix across lead times and shortage penalty regimes.

\begin{figure}[htbp]
\centering
\includegraphics[width=\textwidth]{}
\caption{Portfolio total simulated operating cost preference matrix across lead times and shortage penalty regimes (EXP-019).}
\label{fig:cost_surface}
\end{figure}

Explaining the service-cost trade-off shown in Table~\ref{tab:shortage_sensitivity} and Figure~\ref{fig:cost_surface}: whenever unit shortage penalties exceed $\$2.00/\text{unit}$ (representing high-margin products or steep customer defection costs), conformal replenishment achieves \textbf{lower total operating costs across multi-day horizons}, generating net annual savings of \textbf{--\$515,118} at $L=3, p=\$2.00$, \textbf{--\$2,371,534} at $L=3, p=\$4.00$, and \textbf{--\$1,099,802} at $L=5, p=\$4.00$.

\subsection{Direct Synthesis of Research Questions}
\label{sec:res_rq_synthesis}
Table~\ref{tab:rq_findings} summarizes how the empirical findings resolve RQs 1-4.

\begin{table}[htbp]
\centering
\small
\caption{Systematic Summary of Research Questions, Methodological Implementations, and Empirical Findings.}
\label{tab:rq_findings}
\begin{tabular}{p{1.2cm}p{4.2cm}p{4.5cm}p{4.8cm}}
\toprule
\textbf{RQ} & \textbf{Focus Area} & \textbf{Methodological Implementation} & \textbf{Key Empirical Finding} \\
\midrule
\textbf{RQ1} & Decision-Time Multi-Horizon Forecasting Stability & Seven direct XGBoost models ($M_1\dots M_7$) with $\tau \ge 7$ days constraint & Stable accuracy across horizons (Val MAE 6.0388, RMSE 7.8638); passed 10/10 automated leakage checks. \\
\textbf{RQ2} & Conformal Calibration Under Temporal Distribution Shift & Rolling Forecast-Conditional Conformal Calibration ($W=100, K^*=100$) & Diagnosed $3.55\times\dots4.30\times$ residual std dev expansion; coverage dropped to $18.9\%-32.0\%$; Winkler improved by 5.2\% ($p < 10^{-27}$). \\
\textbf{RQ3} & Replenishment Service and Shortage Mitigation & Dynamic safety stock ($SS = q \cdot \hat{D}$) in discrete-event continuous review & Portfolio fill rate increased by $+7.29$ to $+10.64$ pp ($p < 10^{-75}$); annual lost sales cut by $20.0\%$ to $29.6\%$ (1.14M units saved). \\
\textbf{RQ4} & Economic Cost Trade-Offs and Shortage Sensitivity & Financial accounting ($K_{\text{order}}=\$10, h=\$0.20, p \in [\$0.60, \$4.00]$) & Net cost reduced by $-10.83\%$ at $L=1$ ($-\$790.8\text{k}$); mapped multi-day cost-preference frontier ($p \ge \$2.00$). \\
\bottomrule
\end{tabular}
\end{table}

As summarized in Table~\ref{tab:rq_findings}, decoupled direct forecasting maintains temporal decision-time safety (RQ1); empirical coverage degrades under nonstationary residual-scale expansion while relative scaling preserves interval efficiency (RQ2); uncertainty-derived safety stock delivers substantial, statistically significant fill-rate and shortage improvements (RQ3); and economic cost preference is governed by lead time and shortage penalties, establishing a clear crossover at $p \ge \$2.00$ (RQ4).

% ==============================================================================
% SECTION 5: CONCLUSION, LIMITATIONS, AND FUTURE WORK
% ==============================================================================
\section{Conclusion, Limitations, and Future Work}
\label{sec:conclusion}

\subsection{Summary of the Study}
\label{sec:concl_summary}
This study presented an end-to-end empirical framework bridging leakage-controlled multi-horizon forecasting, rolling forecast-conditional conformal prediction, and continuous-review inventory simulation across 500 store-item time series. The investigation demonstrated that while high point-forecasting accuracy can be maintained under decision-time constraints ($\text{MAE} \approx 6.04$), temporal distribution shift causes severe residual-scale expansion ($3.55\times\dots4.30\times$), breaking nominal conformal coverage guarantees. Nevertheless, translating relative forecast-conditional uncertainty into dynamic safety stock produced major operational service improvements ($+10.64$ pp fill rate at $L=1$, $1.14$ million lost-sales units prevented) and established explicit economic trade-offs across lead times and shortage penalty regimes.

\subsection{Managerial Takeaways}
\label{sec:concl_managerial}
The empirical results provide four actionable takeaways for retail supply chain managers: (1) point forecasting accuracy alone is insufficient to protect against inventory stockouts; (2) uncertainty calibration must be continuously monitored for residual scale expansion under nonstationary demand; (3) uncertainty buffers should be evaluated via operational service and financial metrics rather than statistical interval coverage in isolation; and (4) uncertainty-aware safety buffering is economically advantageous for short-cycle replenishment ($L=1$) and high-margin product categories ($p \ge \$2.00$).

\subsection{Grounded Limitations}
\label{sec:concl_limitations}
To maintain scientific rigor, several study limitations are explicitly stated:
\begin{enumerate}
    \item \textbf{Single Benchmark Dataset:} Evaluation was conducted on 500 series from 10 stores within a single public benchmark dataset.
    \item \textbf{Controlled Simulation Environment:} The discrete-event simulation assumes unconstrained supplier capacity, deterministic delivery lead times, and lost-sales dynamics without backlogging or product substitution.
    \item \textbf{Model Backbone Scope:} The forecasting layer utilized GBDT regressors (XGBoost); deep neural sequence models (TFT, DeepAR) were not benchmarked within the multi-horizon inventory simulation.
    \item \textbf{Conformal Breakdown Under Shift:} The uncertainty framework does not provide guaranteed nominal coverage under unmodeled temporal distribution shift.
    \item \textbf{Cost Function Linearity:} Holding and shortage costs were modeled as fixed linear parameters without non-linear volume discounts or contractual minimum order quantities.
\end{enumerate}

\subsection{Directions for Future Research}
\label{sec:concl_future}
Future research directions include: (1) developing adaptive conformal inference layers with explicit volatility tracking \cite{gibbs2021adaptive}; (2) benchmarking against deep probabilistic and foundation time-series models; (3) extending the simulation engine to multi-echelon networks with stochastic lead times; and (4) evaluating empirical cross-retailer datasets.

% ==============================================================================
% BACK MATTER: DECLARATIONS
% ==============================================================================
\section*{Declarations}

\subsection*{Declaration of Prior Conference Publication and Substantial Extension}
\addcontentsline{toc}{subsection}{Declaration of Prior Conference Publication and Substantial Extension}

\noindent \textbf{Prior Conference Publication:} \\
An earlier preliminary version of this research appeared as:
\begin{quote}
\small
R. Mangaiyarkarasi et al., ``Explainable Intelligent Decision Support System (E-IDSS) for Retail Sales Forecasting and Inventory Optimization,'' in \emph{Proceedings of the IEEE International Conference on Intelligent Decision Support Systems}, 2026, pp. 1-8 \cite{e_idss_ieee}.
\end{quote}

\noindent \textbf{Substantial Methodological and Empirical Extensions in this Journal Manuscript:} \\
The present manuscript extends the preliminary conference study by retaining the underlying retail demand forecasting foundation (the Store Item Demand Forecasting dataset spanning 500 store-item series) while introducing a fully realized, leakage-controlled, distribution-free uncertainty quantification and discrete-event inventory decision-support framework. Specifically, this journal article provides the following substantial extensions beyond the conference baseline:

\begin{enumerate}
    \item \textbf{Decision-Time-Safe Feature Protocol:} Replaces standard unconstrained lag features with a strictly causal temporal feature protocol enforcing a minimum lag constraint of $\tau \ge 7$ days, passing a 10-check automated leakage audit suite.
    \item \textbf{Decoupled Direct Multi-Horizon Architecture:} Replaces the single-step forecasting model with seven independently trained global direct XGBoost regressors ($M_1, \dots, M_7$), eliminating recursive error propagation across lead-time horizons.
    \item \textbf{Out-of-Sample Residual Calibration Pool:} Replaces heuristic in-sample variance estimates with a rolling-origin temporal cross-validation residual calibration pool ($N = 608,165$ residuals) generated without forward test access.
    \item \textbf{Rolling Forecast-Conditional Conformal Calibration:} Introduces a non-parametric conformal prediction layer that dynamically scales nonconformity scores by point-forecast magnitudes ($s = |e| / (\hat{D} + \epsilon)$), benchmarked against raw and scale-normalized conformal variants.
    \item \textbf{Distribution Shift and Scale Expansion Diagnostics:} Provides a formal two-sample location-scale decomposition and Kolmogorov-Smirnov hypothesis testing ($D = 0.170\dots0.197, p < 10^{-42}$) that empirically diagnoses a $3.55\times$ to $4.30\times$ residual standard-deviation expansion and quarterly seasonal error swings on the 2017 test holdout.
    \item \textbf{Calibration Robustness and Neighborhood Sweeps:} Conducts an extensive validation sensitivity sweep across calibration neighborhood sizes $K \in [10, 100]$, identifying a broad stability plateau ($K \in [30, 100]$) and freezing $K^* = 100$ under a pre-registered finite-sample sizing rule.
    \item \textbf{Discrete-Event Inventory Replenishment Simulator:} Translates upper conformal prediction bounds into dynamic safety stock buffers ($SS_{t,L} = q_{t,L} \cdot \hat{D}_{t,L}$) within a continuous-review order-up-to policy ($ROP = \hat{D} + SS$) simulated over 365 daily discrete events under lost-sales dynamics.
    \item \textbf{Multi-Horizon Operational and Cost Evaluation:} Evaluates inventory performance across four supplier delivery lead times ($L \in \{1, 3, 5, 7\}$ days) and four shortage penalty regimes ($p \in \{\$0.60, \$1.00, \$2.00, \$4.00\}$), revealing explicit service--cost trade-offs and identifying the economic penalty crossover threshold ($p \ge \$2.00$) where uncertainty buffering lowers total simulated operating costs.
    \item \textbf{Inferential Statistical and Bootstrap Rigor:} Incorporates non-parametric paired Wilcoxon signed-rank hypothesis tests ($p < 10^{-75}$ for service gains, $p = 6.65 \times 10^{-55}$ for $L=1$ cost reductions) and 1,000-iteration paired bootstrap confidence intervals.
\end{enumerate}

\subsection*{Declaration of Competing Interest}
The authors declare that they have no known competing financial interests or personal relationships that could have appeared to influence the work reported in this paper.

\subsection*{Funding Statement}
This research received no specific grant from any funding agency in the public, commercial, or not-for-profit sectors.

\subsection*{Data and Code Availability Statement}
The Store Item Demand Forecasting dataset is publicly accessible via standard machine learning repositories. Code and experimental pipelines are archived under institutional guidelines.

\subsection*{CRediT Author Statement}
\textbf{R. Mangaiyarkarasi:} Conceptualization, Methodology, Software, Validation, Formal Analysis, Investigation, Writing: Original Draft, Writing: Review \& Editing. \textbf{Research Collaborators:} Methodology, Supervision, Project Administration, Writing: Review \& Editing.

% ==============================================================================
% APPENDICES
% ==============================================================================
\clearpage
\appendix
\section{Mathematical Formulation and Notation Summary}
\label{app:notation}

Table~\ref{tab:app_notation} summarizes the primary mathematical symbols, dimensions, and operational definitions utilized across the manuscript.

\begin{table}[htbp]
\centering
\small
\caption{Summary of Primary Mathematical Notation.}
\label{tab:app_notation}
\begin{tabular}{llp{8cm}}
\toprule
\textbf{Symbol} & \textbf{Domain} & \textbf{Operational Description} \\
\midrule
$i \in \{1,\dots,M\}$ & Integer ($M=500$) & Index of individual store-item time series. \\
$t \in \{1,\dots,T\}$ & Integer & Daily forecast origin / decision timestamp. \\
$h \in \{1,\dots,H\}$ & Integer ($H=7$) & Future forecasting horizon in days ($h \in \{1, \dots, 7\}$). \\
$L \in \{1, 3, 5, 7\}$ & Integer (Days) & Supplier order delivery lead time. \\
$y_{i,t+h}$ & $\mathbb{R}_{\ge 0}$ & Realized retail demand / sales quantity for series $i$ at day $t+h$. \\
$\hat{y}_{i,t+h}$ & $\mathbb{R}_{\ge 0}$ & Point forecast generated by decoupled direct model $M_h$ at origin $t$. \\
$\hat{D}_{i,t,L}$ & $\mathbb{R}_{\ge 0}$ & Cumulative expected lead-time demand: $\hat{D}_{i,t,L} = \sum_{k=1}^L \hat{y}_{i,t+k}$. \\
$\mathcal{F}_{i,t}$ & $\mathbb{R}^{45}$ & Decision-time-safe feature vector constructed from history up to $t-7$. \\
$\tau \ge 7$ & Integer (Days) & Minimum lag constraint enforcing temporal decision-time isolation. \\
$s_{i,r,L}$ & $\mathbb{R}_{\ge 0}$ & Relative forecast-conditional nonconformity score: $|e| / (\hat{D} + \epsilon)$. \\
$q_{i,t,L}$ & $\mathbb{R}_{\ge 0}$ & Empirical $(1-\alpha)$ quantile of nonconformity scores over window $W=100$. \\
$SS_{i,t,L}$ & $\mathbb{R}_{\ge 0}$ & Uncertainty-derived dynamic safety stock buffer: $SS_{i,t,L} = q_{i,t,L} \cdot \hat{D}_{i,t,L}$. \\
$ROP_{i,t,L}$ & $\mathbb{R}_{\ge 0}$ & Continuous-review reorder point: $ROP_{i,t,L} = \hat{D}_{i,t,L} + SS_{i,t,L}$. \\
$S_{i,t,L}$ & $\mathbb{R}_{\ge 0}$ & Target order-up-to inventory level: $S_{i,t,L} = ROP_{i,t,L} + Q_{\text{base}}$. \\
$IP_{i,t}$ & $\mathbb{R}_{\ge 0}$ & Inventory position: $IP_{i,t} = OH_{i,t} + OO_{i,t}$. \\
$OH_{i,t}, OO_{i,t}$ & $\mathbb{R}_{\ge 0}$ & On-hand inventory and on-order (in-transit) inventory at origin $t$. \\
$K_{\text{order}}$ & $\$$ & Fixed order placement transaction cost (\$10.00/order). \\
$h$ & $\$$/unit/day & Unit inventory holding cost (\$0.20/unit/day). \\
$p$ & $\$$/unit & Unit shortage / lost-sales penalty ($p \in \{\$0.60, \$1.00, \$2.00, \$4.00\}$). \\
$TC$ & $\$$ & Total simulated operating cost: $TC = C_{\text{order}} + C_{\text{holding}} + C_{\text{shortage}}$. \\
\bottomrule
\end{tabular}
\end{table}

\section{Complete Feature Engineering Protocol and Category Breakdown}
\label{app:features}

To guarantee complete decision-time safety, all 45 feature dimensions are constructed using historical observations realized strictly at or before $t-7$ days. Table~\ref{tab:app_features} provides the complete category breakdown.

\begin{table}[htbp]
\centering
\small
\caption{Decision-Time-Safe Feature Space Breakdown (45 Dimensions).}
\label{tab:app_features}
\begin{tabular}{lp{9cm}c}
\toprule
\textbf{Feature Category} & \textbf{Specific Mathematical Transformations} & \textbf{Dimension Count} \\
\midrule
Autoregressive Lags & Lags $\tau \in \{7, 8, 9, 10, 11, 12, 13, 14, 21, 28, 35, 42, 60, 90, 180, 365\}$ & 16 \\
Rolling Statistics & Trailing means over windows $\{7, 14, 28, 60, 90, 180\}$ lagged by 7 days & 6 \\
Rolling Volatility & Trailing standard deviations over windows $\{7, 14, 28, 60, 90\}$ lagged by 7 & 5 \\
Rolling Extremes & Trailing minimums and maximums over $\{7, 28\}$ days lagged by 7 days & 4 \\
Exponential Smoothing & Trailing EMAs with smoothing parameters $\alpha \in \{0.1, 0.3, 0.5\}$ lagged by 7 & 3 \\
Calendar Cycles & Sine and cosine encodings for Day of Week, Day of Month, Month of Year & 6 \\
Binary Indicators & Weekend indicator, Month-start indicator, Month-end indicator & 3 \\
Categorical Encodings & Out-of-fold target encoding for Store ID and Item ID & 2 \\
\midrule
\textbf{Total Dimensionality} & \textbf{All features strictly satisfy $\tau \ge 7$ days constraint} & \textbf{45} \\
\bottomrule
\end{tabular}
\end{table}

\section{Discrete-Event Continuous-Review Inventory Simulation Algorithm}
\label{app:algorithm}

Algorithm~\ref{alg:inventory_simulation} details the daily execution cycle executed for each store-item series $i$ over the 365 days of the out-of-sample evaluation period.

\begin{algorithm}[htbp]
\small
\caption{Daily Discrete-Event Replenishment Simulation with Lost-Sales Accounting}
\label{alg:inventory_simulation}
\begin{algorithmic}[1]
\REQUIRE Series demand $\{y_{i,t}\}_{t=1}^T$, Lead time $L$, Unit costs $(K_{\text{order}}, h, p)$, Initial stock $OH_{i,0} = L \cdot \bar{D}_{\text{train}}$, $OO_{i,0} = 0$.
\FOR{day $t = 1$ \TO $T$}
    \STATE \textbf{1. Receive Inbound Shipments:} $OH_{i,t} \leftarrow OH_{i,t} + \text{InboundShipments}(t)$
    \STATE \textbf{2. Update In-Transit Pipeline:} $OO_{i,t} \leftarrow OO_{i,t} - \text{InboundShipments}(t)$
    \STATE \textbf{3. Calculate Inventory Position:} $IP_{i,t} \leftarrow OH_{i,t} + OO_{i,t}$
    \STATE \textbf{4. Generate Point Forecasts:} $\hat{D}_{i,t,L} \leftarrow \sum_{k=1}^L \hat{y}_{i,t+k}$
    \STATE \textbf{5. Compute Dynamic Safety Stock:} 
    \IF{Policy is Conformal}
        \STATE $q_{i,t,L} \leftarrow \text{Quantile}_{1-\alpha}\left( \{s_{i,r,L}\}_{r \in \mathcal{W}_t} \right)$
        \STATE $SS_{i,t,L} \leftarrow q_{i,t,L} \cdot \hat{D}_{i,t,L}$
    \ELSE
        \STATE $SS_{i,t,L} \leftarrow 0$
    \ENDIF
    \STATE \textbf{6. Evaluate Reorder Condition:} $ROP_{i,t,L} \leftarrow \hat{D}_{i,t,L} + SS_{i,t,L}$
    \IF{$IP_{i,t} \le ROP_{i,t,L}$}
        \STATE $Q_{i,t} \leftarrow \max\left(0, \lceil (ROP_{i,t,L} + Q_{\text{base}}) - IP_{i,t} \rceil\right)$
        \STATE Schedule arrival of $Q_{i,t}$ units at day $t+L$
        \STATE $OO_{i,t} \leftarrow OO_{i,t} + Q_{i,t}$
        \STATE Accumulate order placement transaction cost: $C_{\text{order}} \leftarrow C_{\text{order}} + K_{\text{order}}$
    \ENDIF
    \STATE \textbf{7. Customer Demand Realization:} Realized demand $y_{i,t}$ arrives
    \STATE \textbf{8. Order Fulfillment \& Lost-Sales Accounting:}
    \STATE $\text{Sales}_{i,t} \leftarrow \min(OH_{i,t}, y_{i,t})$
    \STATE $\text{LostSales}_{i,t} \leftarrow y_{i,t} - \text{Sales}_{i,t}$
    \STATE $OH_{i,t+1} \leftarrow OH_{i,t} - \text{Sales}_{i,t}$
    \STATE Accumulate shortage penalty: $C_{\text{shortage}} \leftarrow C_{\text{shortage}} + p \cdot \text{LostSales}_{i,t}$
    \STATE Accumulate inventory holding cost: $C_{\text{holding}} \leftarrow C_{\text{holding}} + h \cdot OH_{i,t+1}$
\ENDFOR
\ENSURE Cumulative costs $(C_{\text{order}}, C_{\text{holding}}, C_{\text{shortage}})$, Total Cost $TC$, Fill Rate $\text{FR}_i$, Stockouts.
\end{algorithmic}
\end{algorithm}

\section{Computational Profiling and Hardware Execution Environment}
\label{app:computational_profiling}

Table~\ref{tab:app_computational} details the runtime benchmarks and computational complexity across the stages of the pipeline.

\begin{table}[htbp]
\centering
\small
\caption{Computational Complexity and Pipeline Runtime Benchmarks.}
\label{tab:app_computational}
\begin{tabular}{lcp{8.5cm}}
\toprule
\textbf{Pipeline Stage} & \textbf{Wall-Clock Time} & \textbf{Complexity / Hardware Specifications} \\
\midrule
Feature Engineering (45 features $\times$ 913k obs) & 4.2 Minutes & $\mathcal{O}(N \log N)$ pandas vectorized operations. \\
XGBoost Model Training ($M_1 \dots M_7$, 500 trees) & 24.2 Minutes & $\mathcal{O}(M \cdot d \cdot N_{\text{trees}})$ parallelized over 8 CPU cores. \\
OOS Residual Pooling ($N = 608,165$ residuals) & 8.5 Minutes & 5-fold TimeSeriesSplit cross-validation per model. \\
Rolling Conformal Calibration ($W=100$, 500 SKUs) & 1.8 Minutes & Dynamic quantile extraction across sliding pools. \\
365-Day Inventory Simulation (500 SKUs $\times$ 4 Lead Times) & 3.4 Minutes & Daily event-driven state transitions with lost sales. \\
\midrule
\textbf{Total End-to-End Pipeline Execution} & \textbf{42.1 Minutes} & \textbf{Standard Workstation (x86\_64 CPU, 32 GB RAM)} \\
\bottomrule
\end{tabular}
\end{table}

% ==============================================================================
% EMBEDDED BIBLIOGRAPHY
% ==============================================================================
\clearpage
\begin{thebibliography}{99}
\addcontentsline{toc}{section}{References}

\bibitem{tiwari2025agentic}
S.~Tiwari and R.~Prasad, ``Agentic AI workflows in modern retail supply chain demand forecasting,'' in \emph{Proc. Int. Conf. Artif. Intell. Mach. Vision (AIMV)}, 2025, pp. 230-244.

\bibitem{mitra2022comparative}
S.~Mitra, A.~Chattopadhyay, and D.~Banerjee, ``A comparative study of machine learning approaches for store-item demand forecasting,'' \emph{Oper. Res. Forum}, vol.~3, no.~2, p.~28, 2022.

\bibitem{venkiteela2025machine}
R.~Venkiteela, ``Machine learning architectures for store-item demand estimation under complex seasonality,'' \emph{Int. J. Comput. Eng. Sci.}, vol.~15, no.~1, pp. 88-102, 2025.

\bibitem{mustapha2025forecasting}
A.~Mustapha and M.~Sithole, ``Demand forecasting in retail supply chain management: A comprehensive empirical study on global machine learning models,'' \emph{Afr. J. Stat. Appl. Sci.}, vol.~14, no.~2, pp. 89-110, 2025.

\bibitem{spanos2022application}
A.~Spanos, C.~Vlachokostas, and A.~Karagiannidis, ``Application of artificial intelligence techniques for short-term retail demand forecasting: A review,'' \emph{Comput. Oper. Res.}, vol.~143, p. 105781, 2022.

\bibitem{malik2024sales}
P.~Malik and S.~Kaur, ``Sales forecasting models under non-stationary retail demand: An empirical benchmark,'' \emph{J. Retail. Consum. Serv.}, vol.~78, p. 103720, 2024.

\bibitem{gneiting2007strictly}
T.~Gneiting and A.~E. Raftery, ``Strictly proper scoring rules, prediction, and estimation,'' \emph{J. Amer. Stat. Assoc.}, vol. 102, no. 477, pp. 359-378, 2007.

\bibitem{ban2019big}
G.-Y. Ban and C.~Rudin, ``The big data newsvendor: Practical insights from machine learning,'' \emph{Oper. Res.}, vol.~67, no.~1, pp. 90-108, 2019.

\bibitem{bertsimas2020predictive}
D.~Bertsimas and N.~Kallus, ``From predictive to prescriptive analytics,'' \emph{Manag. Sci.}, vol.~66, no.~3, pp. 1025-1044, 2020.

\bibitem{yeesin2024application}
T.~Yeesin and K.~Srivawat, ``Application of hybrid statistical and machine learning models for retail demand forecasting,'' \emph{Int. J. Logist. Syst. Manag.}, vol.~47, no.~3, pp. 312-330, 2024.

\bibitem{makridakis2022m5}
S.~Makridakis, E.~Spiliotis, and V.~Assimakopoulos, ``The M5 accuracy competition: Results, findings, and conclusions,'' \emph{Int. J. Forecast.}, vol.~38, no.~4, pp. 1346-1364, 2022.

\bibitem{chen2016xgboost}
T.~Chen and C.~Guestrin, ``XGBoost: A scalable tree boosting system,'' in \emph{Proc. 22nd ACM SIGKDD Int. Conf. Knowl. Discovery Data Mining (KDD)}, 2016, pp. 785-794.

\bibitem{hobor2026comparative}
A.~Hobor, ``A comparative benchmark of gradient boosting and temporal neural networks for multi-horizon SKU demand forecasting,'' \emph{Expert Syst. Appl.}, vol. 268, p. 126100, 2026.

\bibitem{alghurabli2025ai}
M.~Alghurabli and A.~Al-Zoubi, ``AI-driven global gradient boosting architectures for multi-echelon retail inventory forecasting,'' \emph{Int. J. Inf. Manag. Data Insights}, vol.~5, no.~1, p. 100210, 2025.

\bibitem{tiwari2026explainable}
S.~Tiwari, ``Explainable multi-horizon machine learning forecasting under decision-time availability constraints,'' \emph{IEEE Trans. Eng. Manag.}, vol.~73, pp. 412-426, 2026.

\bibitem{e_idss_ieee}
R.~Mangaiyarkarasi, S.~Mohankumar, S.~Kiruthika, and N.~Raghapriya, ``Explainable Intelligent Decision Support System (E-IDSS) for Retail Sales Forecasting and Inventory Optimization,'' in \emph{Proc. IEEE Int. Conf. Intell. Decision Support Syst. (E-IDSS)}, 2026, pp. 1-8.

\bibitem{lundberg2017unified}
S.~M. Lundberg and S.-I. Lee, ``A unified approach to interpreting model predictions,'' in \emph{Adv. Neural Inf. Process. Syst. (NeurIPS)}, vol.~30, 2017, pp. 4765-4774.

\bibitem{lin2026critical}
C.~Lin and H.~Wang, ``Critical evaluation of explainable AI methods in enterprise supply chain decision support systems,'' \emph{Decis. Support Syst.}, vol. 176, p. 114050, 2026.

\bibitem{narlawar2025explainable}
P.~Narlawar and S.~Gupta, ``Explainable machine learning framework for retail demand forecasting and stockout prevention,'' \emph{Comput. Ind. Eng.}, vol. 189, p. 109920, 2025.

\bibitem{chang2025explainable}
Y.~Chang and L.~Liu, ``Explainable AI for inventory management: Bridging the trust gap in algorithmic replenishment,'' \emph{Int. J. Prod. Econ.}, vol. 270, p. 109180, 2025.

\bibitem{vovk2005algorithmic}
V.~Vovk, A.~Gammerman, and G.~Shafer, \emph{Algorithmic Learning in a Random World}. New York, NY, USA: Springer, 2005.

\bibitem{angelopoulos2021gentle}
A.~N. Angelopoulos and S.~Bates, ``A gentle introduction to conformal prediction and distribution-free uncertainty quantification,'' \emph{arXiv preprint arXiv:2107.07561}, 2021.

\bibitem{silver1998inventory}
E.~A. Silver, D.~F. Pyke, and R.~Peterson, \emph{Inventory Management and Production Planning and Scheduling}, 3rd ed. New York, NY, USA: John Wiley \& Sons, 1998.

\bibitem{stankeviciute2021conformal}
K.~Stankeviciute, A.~M. Alaa, and M.~van der Schaar, ``Conformal time-series forecasting,'' in \emph{Proc. Adv. Neural Inf. Process. Syst. (NeurIPS)}, vol.~34, 2021, pp. 6216-6228.

\bibitem{huber2019integrating}
J.~Huber and H.~Stuckenschmidt, ``Integrating machine learning and optimization for inventory control under non-stationary demand,'' \emph{Eur. J. Oper. Res.}, vol. 278, no. 3, pp. 883-896, 2019.

\bibitem{barber2023conformal}
R.~F. Barber, E.~J. Cand{\`e}s, A.~Ramdas, and R.~J. Tibshirani, ``Conformal prediction beyond exchangeability,'' \emph{Ann. Stat.}, vol.~51, no.~2, pp. 816-845, 2023.

\bibitem{tibshirani2019conformal}
R.~J. Tibshirani, R.~F. Barber, E.~J. Cand{\`e}s, and A.~Ramdas, ``Conformal prediction under covariate shift,'' in \emph{Adv. Neural Inf. Process. Syst. (NeurIPS)}, vol.~32, 2019, pp. 2530-2540.

\bibitem{syntetos2016supply}
A.~A. Syntetos, Z.~Babai, J.~E. Boylan, S.~Kolassa, and K.~Nikolopoulos, ``Supply chain forecasting: Theory, practice, their gap and the future,'' \emph{Eur. J. Oper. Res.}, vol. 252, no. 1, pp. 1-26, 2016.

\bibitem{axsaeter2015inventory}
S.~Axs{\"a}ter, \emph{Inventory Control}, 3rd ed. Cham, Switzerland: Springer International Publishing, 2015.

\bibitem{snyder2019fundamentals}
L.~V. Snyder and Z.-J.~M. Shen, \emph{Fundamentals of Supply Chain Theory}, 2nd ed. Hoboken, NJ, USA: John Wiley \& Sons, 2019.

\bibitem{gibbs2021adaptive}
I.~Gibbs and E.~Cand{\`e}s, ``Adaptive conformal inference under distribution shift,'' in \emph{Adv. Neural Inf. Process. Syst. (NeurIPS)}, vol.~34, 2021, pp. 1660-1672.

\bibitem{chaganti2026}
R.~Chaganti and V.~Kumar, ``Human-in-the-loop explainable AI for enterprise demand planning,'' \emph{J. Bus. Logist.}, vol.~47, no.~1, pp. 55-72, 2026.

\bibitem{madabhushini2025}
S.~Madabhushini, ``Game-theoretic feature attribution in high-dimensional retail forecasting,'' \emph{IEEE Trans. Knowl. Data Eng.}, vol.~37, no.~4, pp. 1820-1834, 2025.

\bibitem{swami2020}
A.~Swami and P.~Kaur, ``Evaluating tree-based ensemble methods for multi-store retail sales estimation,'' \emph{Int. J. Inf. Technol.}, vol.~12, no.~3, pp. 789-798, 2020.

\bibitem{salinas2020deepar}
D.~Salinas, V.~Flunkert, H.~Januschowski, and C.~Gasthaus, ``DeepAR: Probabilistic forecasting with autoregressive recurrent networks,'' \emph{Int. J. Forecast.}, vol.~36, no.~3, pp. 1181-1191, 2020.

\end{thebibliography}

\end{document}
